\documentclass[a4paper,12pt]{article}
\usepackage[T1]{fontenc}
\usepackage[utf8]{inputenc}
\usepackage[ngerman]{babel}
\usepackage[a4paper,margin=2.2cm]{geometry}
\usepackage{amsmath,amssymb,mathtools}
\usepackage{siunitx}
\usepackage{physics}
\usepackage{bm}
\usepackage{booktabs}
\usepackage{array}
\usepackage{tabularx}
\usepackage{tikz}
\usepackage{pgfplots}
\usepackage{xcolor}
\usepackage{enumitem}
\usepackage{hyperref}
\usepackage{fancyhdr}
\usepackage{titlesec}
\pgfplotsset{compat=1.18}
\usetikzlibrary{arrows.meta,calc,decorations.pathreplacing,positioning}

\hypersetup{colorlinks=true,linkcolor=black,urlcolor=blue}
\setlength{\parindent}{0pt}
\setlength{\parskip}{0.5em}
\setlength{\headheight}{15pt}
\pagestyle{fancy}
\fancyhf{}
\lhead{Physik I}
\rhead{Kinematik eines Massepunkts}
\cfoot{\thepage}

\definecolor{accent}{RGB}{35,70,130}
\definecolor{lightaccent}{RGB}{232,238,248}

\titleformat{\section}{\large\bfseries\color{accent}}{\thesection}{0.6em}{}
\titleformat{\subsection}{\normalsize\bfseries\color{accent}}{\thesubsection}{0.6em}{}
\titleformat{\subsubsection}{\normalsize\bfseries}{\thesubsubsection}{0.6em}{}

\newcommand{\note}[1]{%
  \begin{center}
  \fcolorbox{accent}{lightaccent}{%
    \begin{minipage}{0.94\linewidth}
    \textbf{Ergänzung aus den Vorlesungsnotizen:} #1
    \end{minipage}}
  \end{center}}

\begin{document}

\begin{titlepage}
\centering
{\Huge \textbf{Physik I}}\\[0.4cm]
{\LARGE Kinematik eines Massepunkts}\\[0.8cm]
{\large Zusammengeführte Mitschrift aus PDF-Folie und handschriftlichen Vorlesungsnotizen}\\[1.2cm]
\vfill
\begin{minipage}{0.9\linewidth}
\small
Dieses Dokument fasst die Inhalte der hochgeladenen PDF mit den während der Vorlesung ergänzten handschriftlichen Notizen in einer einheitlichen \LaTeX-Datei zusammen. Inhaltlich geht es um Ortsvektor, Geschwindigkeit, Beschleunigung, gleichmäßig beschleunigte Bewegung, freien Fall und zusammengesetzte Bewegungen.
\end{minipage}
\vfill
{\large \today}
\end{titlepage}

\tableofcontents
\newpage

\section{Kinematik eines Massepunkts}

Ein \textbf{Massepunkt} ist ein materieller Punkt ohne Ausdehnung. In der Kinematik wird nur die Bewegung beschrieben, noch nicht ihre Ursache.

\subsection{Bewegungsgleichungen}
\subsubsection{Ort, Geschwindigkeit, Beschleunigung}

Ort und Bewegung sind relativ; sie hängen vom gewählten \textbf{Bezugssystem} ab.

\note{In den handschriftlichen Notizen steht ausdrücklich, dass man zuerst ein geeignetes räumliches Koordinatensystem festlegt und dass Ort und Bewegung davon abhängen, \emph{welches} Koordinatensystem gewählt wird.}

\paragraph{(i) Räumliches Koordinatensystem festlegen}
Wir wählen ein kartesisches Koordinatensystem mit Ursprung $O=(0,0,0)$ und den Achsen $x$, $y$, $z$.

\begin{center}
\begin{tikzpicture}[scale=1.25,>=Latex]
  \draw[->,thick] (0,0) -- (3.4,0) node[below] {$x$};
  \draw[->,thick] (0,0) -- (2.2,1.3) node[right] {$y$};
  \draw[->,thick] (0,0) -- (0,2.8) node[left] {$z$};
  \coordinate (A) at (1.4,1.55);
  \coordinate (B) at (2.5,1.95);
  \draw[->,very thick,accent] (0,0) -- (A) node[above left] {$\vec s_0(t_0)$};
  \draw[->,very thick,accent] (0,0) -- (B) node[above right] {$\vec s(t_1)$};
  \draw[densely dashed] (A) -- (B);
  \fill (A) circle (1.5pt);
  \fill (B) circle (1.5pt);
  \node[below left] at (0,0) {$O(0,0,0)$};
\end{tikzpicture}
\end{center}

Ein Ort wird durch einen \textbf{Ortsvektor} beschrieben.

\paragraph{(ii) Lage durch Koordinaten}
Der Ort eines Punktes wird durch die Achsenabschnitte bzw. Koordinaten $x$, $y$, $z$ festgelegt:
\[
\vec s_0=
\begin{pmatrix}
 x_0\\ y_0\\ z_0
\end{pmatrix}.
\]

\note{Die Notizen ergänzen hier: Diese Schreibweise beschreibt den Ort zu einem bestimmten Zeitpunkt, zum Beispiel zu $t_0$. Bei Bewegung hängt der Ortsvektor vom Zeitpunkt ab.}

Bei einer Bewegung schreibt man daher allgemeiner
\[
\vec s = \vec s(t).
\]

\paragraph{(iii) Geschwindigkeit}
Bei Bewegung ändert sich der Ort mit der Zeit. Die Geschwindigkeit ist die zeitliche Änderung des Ortes.

Für eine eindimensionale Bewegung gilt zunächst als mittlere Geschwindigkeit
\[
 v = \frac{\Delta s}{\Delta t}
 \qquad [m/s]. 
\]
Diese Darstellung ist nur dann exakt als momentane Geschwindigkeit brauchbar, wenn $v$ konstant ist.

Besser ist die Definition über die Ableitung:
\[
 v(t)=\frac{\dd s}{\dd t}.
\]

\note{In den Notizen wird dazu ausdrücklich vermerkt: \glqq Zeit-Ableitung liefert die momentane Geschwindigkeit\grqq.}

In drei Dimensionen erhält man den Geschwindigkeitsvektor als Ableitung des Ortsvektors:
\[
 \vec v(t)=\frac{\dd \vec s(t)}{\dd t}=\dot{\vec s}(t)
 =
 \begin{pmatrix}
   v_x\\ v_y\\ v_z
 \end{pmatrix}
 =
 \begin{pmatrix}
   \dfrac{\dd x}{\dd t}\\[0.4em]
   \dfrac{\dd y}{\dd t}\\[0.4em]
   \dfrac{\dd z}{\dd t}
 \end{pmatrix}.
\]

\note{Die Notizen ergänzen, dass man die Koordinaten separat nach der Zeit ableitet; der Punkt über einer Größe steht für die Zeitableitung.}

\paragraph{(iv) Beschleunigung}
Wenn sich die Geschwindigkeit mit der Zeit ändert, liegt eine Beschleunigung vor:
\[
 \vec a(t)=\frac{\dd \vec v(t)}{\dd t}=\dot{\vec v}(t)=\ddot{\vec s}(t).
\]
In Komponenten:
\[
 \vec a(t)=
 \begin{pmatrix}
   \dfrac{\dd v_x}{\dd t}\\[0.4em]
   \dfrac{\dd v_y}{\dd t}\\[0.4em]
   \dfrac{\dd v_z}{\dd t}
 \end{pmatrix}
 =
 \begin{pmatrix}
   \dfrac{\dd^2 x}{\dd t^2}\\[0.4em]
   \dfrac{\dd^2 y}{\dd t^2}\\[0.4em]
   \dfrac{\dd^2 z}{\dd t^2}
 \end{pmatrix}.
\]

Für die eindimensionale Betrachtung gilt:
\[
 a>0 \Rightarrow \text{Beschleunigung},
 \qquad
 a<0 \Rightarrow \text{Bremsen bzw. Verzögerung}.
\]

Zwischen Ort, Geschwindigkeit und Beschleunigung besteht also die Kette
\[
 \vec s(t) \xrightarrow{\ \frac{\dd}{\dd t}\ } \vec v(t)
 \xrightarrow{\ \frac{\dd}{\dd t}\ } \vec a(t),
\]
und umgekehrt erhält man durch Integration wieder $\vec v(t)$ und $\vec s(t)$.

\begin{center}
\begin{tikzpicture}[>=Latex, node distance=4cm]

  \node (s) {$\vec s(t)$};
  \node[right=of s] (v) {$\vec v(t)$};
  \node[right=of v] (a) {$\vec a(t)$};

  % Gerade Pfeile: oben beschriften
  \draw[->,thick] (s) -- node[above=6pt] {Ableitung} (v);
  \draw[->,thick] (v) -- node[above=6pt] {Ableitung} (a);

  % Gebogene Pfeile: unten führen und unten beschriften
  \draw[->,thick] (v) to[out=230,in=310] node[below=10pt] {Aufleitung} (s);
  \draw[->,thick] (a) to[out=230,in=310] node[below=10pt] {Aufleitung} (v);

\end{tikzpicture}
\end{center}



\subsection{Gleichmäßig beschleunigte Bewegung}
Wir betrachten nun eine Bewegung mit konstanter Beschleunigung
\[
 a(t)=a_0=\text{konstant}.
\]

\subsubsection{Messwerte aus dem Versuch}
In den Notizen wurden zu einem Experiment folgende Werte festgehalten:

\begin{center}
\begin{tabular}{>{\centering\arraybackslash}m{2.3cm} >{\centering\arraybackslash}m{2.7cm} >{\centering\arraybackslash}m{2.7cm} >{\centering\arraybackslash}m{2.7cm}}
\toprule
$s$ [\si{\milli\meter}] & 250 & 910 & 2000\\
\midrule
$t$ [\si{\second}] & 0{,}949 & 1{,}962 & 2{,}979\\
\bottomrule
\end{tabular}
\end{center}

Aus
\[
 a_0 = \frac{\dd^2 s}{\dd t^2}= \ddot{\overrightarrow{s}} = \text{konstant}
\]
folgt durch Integration zunächst die Geschwindigkeit:
\[
 v(t)=\int_0^t a_0\,\dd t' = a_0 t + v_0.
\]
Dabei ist $v_0$ die Integrationskonstante, also die Anfangsgeschwindigkeit.

Nochmaliges Integrieren liefert die Weg-Zeit-Funktion:
\[
 s(t)=\int_0^t v(t')\,\dd t'
 =\int_0^t (a_0 t'+v_0)\,\dd t'
 =\frac12 a_0 t^2 + v_0 t + s_0.
\]
Dabei ist $s_0$ die Anfangsposition.

\note{Die handschriftlichen Notizen markieren explizit: $v_0$ und $s_0$ sind Konstanten aus der Integration. Außerdem wird angemerkt, dass die Parabelform von $s(t)$ typisch für konstante Beschleunigung ist.}

Damit gelten die Standardgleichungen der gleichmäßig beschleunigten Bewegung:
\[
 a(t)=a_0,
 \qquad
 v(t)=a_0 t+v_0,
 \qquad
 s(t)=\frac12 a_0 t^2+v_0 t+s_0.
\]

\begin{center}
\begin{tikzpicture}
\begin{axis}[
  width=0.3\linewidth,height=5cm,
  xmin=0,xmax=4,ymin=0,ymax=2,
  axis lines=left,
  xlabel={$t$},ylabel={$a(t)$},
  xtick=\empty,ytick=\empty]
  \addplot[accent,thick,domain=0:4] {1.2};
  \node at (axis cs:0.35,1.42) {$a_0$};
\end{axis}
\end{tikzpicture}
\hfill
\begin{tikzpicture}
\begin{axis}[
  width=0.3\linewidth,height=5cm,
  xmin=0,xmax=4,ymin=0,ymax=4,
  axis lines=left,
  xlabel={$t$},ylabel={$v(t)$},
  xtick=\empty,ytick=\empty]
  \addplot[accent,thick,domain=0:4] {0.7*x + 0.8};
  \node at (axis cs:0.35,0.95) {$v_0$};
\end{axis}
\end{tikzpicture}
\hfill
\begin{tikzpicture}
\begin{axis}[
  width=0.3\linewidth,height=5cm,
  xmin=0,xmax=4,ymin=0,ymax=5,
  axis lines=left,
  xlabel={$t$},ylabel={$s(t)$},
  xtick=\empty,ytick=\empty]
  \addplot[accent,thick,domain=0:4] {0.22*x^2 + 0.28*x + 0.35};
  \node at (axis cs:0.28,0.5) {$s_0$};
\end{axis}
\end{tikzpicture}
\end{center}

\subsection{Freier Fall}
Beim freien Fall ohne Reibung ist die Beschleunigung konstant und gleich der Erdbeschleunigung $g$:
\[
 a=g \approx \SI{9.81}{\meter\per\second\squared}.
\]

Für einen Körper, der aus der Ruhe losgelassen wird, gilt
\[
 v_0=0,
 \qquad s_0=0,
 \qquad v(t)=gt,
 \qquad s(t)=\frac12 gt^2.
\]

Will man die Geschwindigkeit als Funktion der Fallhöhe $h$ bestimmen, setzt man
\[
 h=s(t)=\frac12 gt^2
\qquad\Rightarrow\qquad
 t=\sqrt{\frac{2h}{g}}.
\]
Eingesetzt in $v(t)=gt$ ergibt das
\[
 v(h)=\sqrt{2gh}.
\]

%
%\begin{center}
%\begin{tikzpicture}[>=Latex,scale=1]
%  \draw[thick] (0,0) -- (0,4);
%  \fill (0,3.6) circle (1.8pt);
%  \draw[->,very thick] (0.8,3.5) -- (0.8,0.4);
%  \node[right] at (0.8,2.0) {$a=g$};
%  \draw[decorate,decoration={brace,mirror,amplitude=5pt}] (1.4,3.6) -- (1.4,0.2);
%  \node[right] at (1.55,1.9) {$h(t)$};
%  \node[left] at (0,0) {$0$};
%\end{tikzpicture}
%\end{center}

\paragraph{Beispiel: Stein fällt aus \SI{100}{\meter} Höhe}
\[
 v = \sqrt{2\cdot \SI{9.81}{\frac{m}{s^2}}\cdot \SI{100}{\meter}}
 \approx \SI{44.3}{\frac{m}{s}}
 \approx \SI{159.7}{\frac{km}{h}}.
\]

\note{In den Notizen wurde zusätzlich ein Näherungsversuch mit $t\approx \SI{246}{\milli\second}$ und $s=\SI{30}{\centi\meter}$ festgehalten. Daraus ergibt sich über $g=2s/t^2$ ein Wert von ungefähr $\SI{9.92}{\meter\per\second\squared}$ und damit eine gute Näherung an den Literaturwert $\SI{9.81}{\meter\per\second\squared}$.}

\subsection{Senkrechter Wurf}
Beim senkrechten Wurf ist $v_0\neq 0$ und die Vorzeichenkonvention muss sauber beachtet werden.

\paragraph{Wurf nach unten} Wählt man die positive Richtung nach unten, dann gilt
\[
 v(t)=v_0+gt,
 \qquad
 s(t)=v_0 t+\frac12 gt^2.
\]

\paragraph{Wurf nach oben} Wählt man die positive Richtung nach oben, dann wirkt die Erdbeschleunigung entgegen der Bewegungsrichtung:
\[
 v(t)=v_0-gt,
 \qquad
 s(t)=v_0 t-\frac12 gt^2.
\]

\note{Die Notizen betonen hier ausdrücklich: Auf die Vorzeichen achten.}

\subsection{Zusammengesetzte Bewegung}
Bei zusammengesetzten Bewegungen werden Geschwindigkeitsvektoren vektoriell addiert. Im Beispiel aus der PDF ist eine Vorwärtsbewegung mit einer Bewegung nach oben kombiniert:
\[
 \vec v = \vec v_F + \vec v_B.
\]

Mit
\[
 \vec v_F=
 \begin{pmatrix}
 5\,\si{\meter\per\second}\\0\\0
 \end{pmatrix},
 \qquad
 \vec v_B=
 \begin{pmatrix}
 0\\10\,\si{\meter\per\second}\\0
 \end{pmatrix}
\]
folgt
\[
 \vec v=
 \begin{pmatrix}
 5\\10\\0
 \end{pmatrix}\si{\meter\per\second}.
\]
Der Betrag ist
\[
 \abs{\vec v}=\sqrt{v_x^2+v_y^2+v_z^2}
 =\sqrt{5^2+10^2}\,\si{\meter\per\second}
 \approx \SI{11.2}{\meter\per\second}.
\]

\begin{center}
\begin{tikzpicture}[scale=0.9,>=Latex]
  \draw[->,thick] (0,0) -- (5.2,0) node[below] {$x$};
  \draw[->,thick] (0,0) -- (0,4.8) node[left] {$y$};
  \draw[->,very thick,accent] (0,0) -- (2.2,0) node[midway,below] {$\vec v_F$};
  \draw[->,very thick,accent] (2.2,0) -- (2.2,3.2) node[midway,right] {$\vec v_B$};
  \draw[->,very thick] (0,0) -- (2.2,3.2) node[midway,left] {$\vec v$};
  \draw[densely dashed] (0,3.2) -- (2.2,3.2);
\end{tikzpicture}
\end{center}

\section{Kompakte Zusammenfassung}
\begin{align*}
\vec v(t) &= \dot{\vec s}(t),\\
\vec a(t) &= \dot{\vec v}(t)=\ddot{\vec s}(t),\\
\text{bei } a=a_0\text{ konstant:}\qquad
v(t) &= a_0 t + v_0,\\
 s(t) &= \frac12 a_0 t^2 + v_0 t + s_0,\\
\text{freier Fall aus der Ruhe:}\qquad
v(t) &= gt,\\
 s(t) &= \frac12 gt^2,\\
 v(h) &= \sqrt{2gh}.
\end{align*}

\end{document}
